\documentclass[12pt,a4paper]{article}

% Packages
\usepackage{fontspec}
\usepackage{geometry}
\usepackage{hyperref}
\usepackage{listings}
\usepackage{xcolor}
\usepackage{booktabs}
\usepackage{longtable}
\usepackage{array}
\usepackage{graphicx}
\usepackage{fancyhdr}
\usepackage{titlesec}
\usepackage{enumitem}
\usepackage{mdframed}
\usepackage{tocloft}
\usepackage{microtype}
\usepackage{parskip}
\usepackage{amsmath}
\usepackage{amssymb}
\usepackage{ragged2e}
\usepackage{seqsplit}
\usepackage{etoolbox}

% Page geometry
\geometry{
  a4paper,
  top=2.5cm,
  bottom=2.5cm,
  left=2.5cm,
  right=2.5cm,
  headheight=14pt
}

% Fonts
\setmainfont{DejaVu Serif}
\setmonofont{DejaVu Sans Mono}

% Colors
\definecolor{codegreen}{rgb}{0,0.6,0}
\definecolor{codegray}{rgb}{0.5,0.5,0.5}
\definecolor{codepurple}{rgb}{0.58,0,0.82}
\definecolor{backcolour}{rgb}{0.96,0.96,0.96}
\definecolor{criticalred}{RGB}{220,50,50}
\definecolor{warnorange}{RGB}{230,130,20}
\definecolor{okgreen}{RGB}{30,150,30}
\definecolor{starblue}{RGB}{0,70,140}
\definecolor{tableheadblue}{RGB}{30,80,160}
\definecolor{dirgreen}{RGB}{10,100,40}
\definecolor{prpurple}{RGB}{100,0,160}

% Code listing style
\lstdefinestyle{codestyle}{
    backgroundcolor=\color{backcolour},
    commentstyle=\color{codegreen},
    keywordstyle=\color{starblue}\bfseries,
    numberstyle=\tiny\color{codegray},
    stringstyle=\color{codepurple},
    basicstyle=\ttfamily\footnotesize,
    breakatwhitespace=false,
    breaklines=true,
    captionpos=b,
    keepspaces=true,
    numbers=none,
    showspaces=false,
    showstringspaces=false,
    showtabs=false,
    tabsize=4,
    frame=single,
    rulecolor=\color{gray!40},
    xleftmargin=4pt,
    xrightmargin=4pt,
    aboveskip=6pt,
    belowskip=6pt
}
\lstset{style=codestyle}

% Hyperref config
\hypersetup{
    colorlinks=true,
    linkcolor=starblue,
    urlcolor=starblue,
    citecolor=starblue,
    pdfauthor={STAR Project Team},
    pdftitle={STAR Project: Comprehensive Gap Analysis},
    pdfsubject={Gap analysis and implementation roadmap by directory},
    pdfkeywords={RX72N, ROS2, Gateway, robotics, gap analysis},
}

% Section styling
\titleformat{\section}{\Large\bfseries\color{starblue}}{}{0em}{}[\titlerule]
\titleformat{\subsection}{\large\bfseries\color{starblue!80}}{}{0em}{}
\titleformat{\subsubsection}{\normalsize\bfseries}{}{0em}{}

% Header and footer
\pagestyle{fancy}
\fancyhf{}
\fancyhead[L]{\small\textit{STAR Project -- Gap Analysis \& Roadmap}}
\fancyhead[R]{\small\textit{2026-02-18}}
\fancyfoot[C]{\thepage}
\renewcommand{\headrulewidth}{0.4pt}

% Custom box environments
\newmdenv[
  backgroundcolor=criticalred!10,
  linecolor=criticalred,
  linewidth=2pt,
  roundcorner=4pt,
  skipabove=8pt,
  skipbelow=8pt,
  innerleftmargin=8pt,
  innerrightmargin=8pt,
]{criticalbox}

\newmdenv[
  backgroundcolor=warnorange!10,
  linecolor=warnorange,
  linewidth=1.5pt,
  roundcorner=4pt,
  skipabove=6pt,
  skipbelow=6pt,
  innerleftmargin=8pt,
  innerrightmargin=8pt,
]{warnbox}

\newmdenv[
  backgroundcolor=okgreen!8,
  linecolor=okgreen,
  linewidth=1.5pt,
  roundcorner=4pt,
  skipabove=6pt,
  skipbelow=6pt,
  innerleftmargin=8pt,
  innerrightmargin=8pt,
]{okbox}

\newmdenv[
  backgroundcolor=prpurple!6,
  linecolor=prpurple,
  linewidth=1.5pt,
  roundcorner=4pt,
  skipabove=6pt,
  skipbelow=6pt,
  innerleftmargin=8pt,
  innerrightmargin=8pt,
]{prbox}

% Directory heading style
\newmdenv[
  backgroundcolor=starblue!8,
  linecolor=starblue,
  linewidth=2pt,
  roundcorner=4pt,
  skipabove=12pt,
  skipbelow=6pt,
  innerleftmargin=10pt,
  innerrightmargin=10pt,
  innertopmargin=6pt,
  innerbottommargin=6pt,
]{dirbox}

% Helper commands
\newcommand{\done}{\textcolor{okgreen}{\textbf{DONE}}}
\newcommand{\partialstatus}{\textcolor{warnorange}{\textbf{PARTIAL}}}
\newcommand{\missing}{\textcolor{criticalred}{\textbf{MISSING}}}
\newcommand{\inprogress}[1]{\textcolor{prpurple}{\textbf{PR \##1}}}
% \filepath: render path in green monospace with seqsplit so long paths
% can break at any character boundary when needed in narrow columns.
\newcommand{\filepath}[1]{\textcolor{dirgreen}{\texttt{\seqsplit{#1}}}}

% Table column helpers (p-type with raggedright)
\newcolumntype{L}[1]{>{\RaggedRight\arraybackslash}p{#1}}

% Allow relaxed line-breaking inside all coloured info boxes so long
% \texttt{} tokens don't overflow into the margin.
\AtBeginEnvironment{prbox}{\sloppy\hbadness=10000}
\AtBeginEnvironment{warnbox}{\sloppy\hbadness=10000}
\AtBeginEnvironment{okbox}{\sloppy\hbadness=10000}
\AtBeginEnvironment{criticalbox}{\sloppy\hbadness=10000}
\AtBeginEnvironment{dirbox}{\sloppy\hbadness=10000}

\begin{document}

% ============================================================
% TITLE PAGE
% ============================================================
\begin{titlepage}
  \begin{center}
    \vspace*{2cm}
    {\Huge\bfseries\color{starblue} STAR Project}\\[0.5cm]
    {\Large\bfseries Comprehensive Gap Analysis \& Implementation Roadmap}\\[0.3cm]
    {\large Organized by Codebase Directory}\\[2cm]

    \begin{tabular}{ll}
      \toprule
      \textbf{Date} & 2026-02-18 \\
      \textbf{Branch} & \texttt{feat/gap-analysis-and-planning} \\
      \textbf{Overall Completion} & $\sim$75\% \\
      \textbf{Open PRs} & 9 active (see annotations throughout) \\
      \textbf{Next Milestone} & First robot motion (SPI bridge + UI real data) \\
      \bottomrule
    \end{tabular}

    \vspace{1.5cm}
    \begin{okbox}
      \textbf{Confirmed complete (corrected from initial analysis):}\\
      \vspace{4pt}
      WireMessage dispatcher: \done\ (\filepath{comm\_task.c:1959} -- SetVelocity, E-Stop, PID, RetransmitConfig)\\
      BMS telemetry streaming: \done\ (\filepath{telemetry\_task.c:1833} -- voltage\_v, soc\_percent)
    \end{okbox}

    \vspace{0.6cm}
    \begin{prbox}
      \textbf{Notation:} Items marked \inprogress{NNN} have an active pull request.\\
      \textbf{Scope:} This is a prototype. OTA firmware updates are a future enhancement, not a current requirement.
    \end{prbox}

    \vspace{1cm}
    {\large\textit{STAR (Simultaneous Tracking and Robotics) -- Locked Inc.}}
  \end{center}
\end{titlepage}

% ============================================================
% TABLE OF CONTENTS
% ============================================================
\tableofcontents
\newpage

% ============================================================
% SECTION 1: EXECUTIVE SUMMARY
% ============================================================
\section{Executive Summary}

STAR is a distributed robotics platform with a custom PCB RX72N motor controller, Raspberry Pi 5 control system, ROS2 Jazzy middleware, and a Go gateway service. The system communicates over 10~Mbps SPI using Protocol Buffers with HARQ/FEC error correction.

\subsection{System Architecture}

\begin{center}
\small
\begin{tabular}{L{5cm}L{3cm}L{6cm}}
  \toprule
  \textbf{Directory} & \textbf{Language} & \textbf{Role} \\
  \midrule
  \filepath{e2-studio-star-rx72n-firmware/} & C23, GNURX, ThreadX & Motor controller firmware \\
  \filepath{star-gateway/} & Go & UI/ROS2 bridge, SPI driver, frame protocol \\
  \filepath{star-ros2/} & C++17, ROS2 Jazzy & Robot middleware, SPI bridge, safety monitor \\
  \filepath{star-ui/} & TypeScript, React 19 & Operator interface \\
  \filepath{star-proto/} & Protocol Buffers 3 & Message schemas (10 files, 72 messages) \\
  \filepath{star-rpi5-buildroot/} & Buildroot & Custom Linux image for RPi5 \\
  \filepath{matlab/} & MATLAB & Motor system ID and PID design \\
  \filepath{schematic/} & KiCad & Custom PCB designs (MCU, Power, BMS) \\
  \bottomrule
\end{tabular}
\end{center}

\subsection{Communication Flow}

\begin{center}
\small
\begin{tabular}{L{3cm}@{\hspace{0.4cm}}c@{\hspace{0.4cm}}L{3.5cm}@{\hspace{0.4cm}}c@{\hspace{0.4cm}}L{3.5cm}}
  \toprule
  \textbf{UI} & $\to$ & \textbf{Gateway (Go)} & $\to$ & \textbf{ROS2 (C++)} \\
  WebSocket / REST & & gRPC server on RPi5 & & \texttt{/cmd\_vel}, \texttt{/robot\_status} \\[3pt]
  & & & $\downarrow$ & \\
  & & & \textbf{SPI bridge} & \\
  & & & 10~Mbps SPI & \\
  & & & $\downarrow$ & \\
  & & & \textbf{RX72N} & \\
  & & & C23 + ThreadX & \\
  \bottomrule
\end{tabular}
\end{center}

\subsection{Overall Completion Matrix}

\begin{center}
\small
\begin{longtable}{@{}L{5.9cm}L{2.1cm}L{2.1cm}L{3.6cm}@{}}
  \toprule
  \textbf{Component} & \textbf{Complete} & \textbf{Status} & \textbf{Note / PR} \\
  \midrule
  \endhead

  \multicolumn{4}{l}{\textbf{e2-studio-star-rx72n-firmware/}} \\
  \quad Firmware core (9 ThreadX tasks) & 100\% & \done & \\
  \quad Comm stack (HARQ/FEC/CRC-32) & 100\% & \done & \\
  \quad WireMessage command dispatcher & 100\% & \done & \\
  \quad BMS telemetry streaming & 100\% & \done & Fix: \inprogress{342} \\
  \quad Telemetry USB/SPI failover & 95\% & \inprogress{342} & PR open \\
  \quad ThreadX stack checking & 95\% & \inprogress{343} & PR open \\
  \quad BMS SoC thresholds & 95\% & \inprogress{344} & PR open \\
  \quad Event-flag cleanup & 95\% & \inprogress{345} & PR open \\
  \quad Frame decoder resync & 95\% & \inprogress{341} & PR open \\
  \quad ASCII source cleanup & 95\% & \inprogress{346} & PR open \\
  \quad NVS configuration persistence & 0\% & \missing & Not started \\
  \quad OTA firmware update handler & 0\% & \missing & Future enhancement \\
  \midrule

  \multicolumn{4}{l}{\textbf{star-gateway/}} \\
  \quad Transport layer (SPI/USB CDC) & 100\% & \done & \\
  \quad Frame codec (HARQ/FEC/CRC-32) & 100\% & \done & \\
  \quad MotorControlService gRPC & 90\% & \partialstatus & E-Stop priority queue \\
  \quad TelemetryService gRPC & 90\% & \partialstatus & GetSystemStatus stub \\
  \quad BatteryManagementService gRPC & 100\% & \done & \\
  \quad ConfigurationService gRPC & 85\% & \partialstatus & RetransmitConfig dispatch \\
  \quad FirmwareUpdateService gRPC & 0\% & \missing & Future enhancement \\
  \quad RESET\_ACK auto-response & 95\% & \inprogress{349} & PR open \\
  \quad Virtual RX72N simulator & 95\% & \done & No motor dynamics \\
  \midrule

  \multicolumn{4}{l}{\textbf{star-ros2/}} \\
  \quad star\_gateway\_bridge\_node & 100\% & \done & \\
  \quad star\_spi\_bridge\_node & 85\% & \partialstatus & Never run on hardware \\
  \quad star\_safety\_monitor & 40\% & \partialstatus & 3 of 10 tests skipped \\
  \quad star\_bringup launch files & 5\% & \missing & Empty skeleton \\
  \quad RPLiDAR C1 ROS2 driver & 0\% & \missing & Host-side sensor \\
  \quad SLAM stack (RTAB-Map, EKF) & 0\% & \missing & Requires RPLiDAR + IMU \\
  \midrule

  \multicolumn{4}{l}{\textbf{star-ui/}} \\
  \quad App scaffold + 7 pages & 80\% & \inprogress{351} & Mock data only \\
  \quad Real gRPC / HTTP connection & 0\% & \missing & Mock client only \\
  \quad Emergency Stop button & 0\% & \missing & Safety critical \\
  \quad Tests (unit / E2E) & 0\% & \missing & None yet \\
  \midrule

  \multicolumn{4}{l}{\textbf{star-proto/}} \\
  \quad Proto schemas (72 messages) & 100\% & \done & \\
  \quad Generated code (Go/TS/nanopb/C++) & 100\% & \done & \\
  \quad nanopb motor\_status max\_count & 85\% & \partialstatus & Bug: 2 should be 4 \\
  \midrule

  \multicolumn{4}{l}{\textbf{CI/CD (.github/workflows/)}} \\
  \quad proto.yml & 100\% & \done & \\
  \quad gateway.yml & 100\% & \done & \\
  \quad ros2.yml & 100\% & \done & \\
  \quad ui.yml (npm build) & 0\% & \missing & Not yet created \\
  \quad docs.yml (LaTeX compile) & 0\% & \missing & Not yet created \\
  \bottomrule
\end{longtable}
\end{center}

\subsection{Active Pull Requests Summary}

\begin{center}
\small
\begin{tabular}{L{1.3cm}L{4.8cm}L{5.1cm}L{2.5cm}}
  \toprule
  \textbf{PR} & \textbf{Title} & \textbf{What It Fixes} & \textbf{Status} \\
  \midrule
  \#341 & Frame decoder resync & Bounded sync-word recovery & Tests passing \\
  \#342 & USB$\to$SPI telemetry failover & Telemetry was USB-only in prod & Tests passing \\
  \#343 & ThreadX stack checking & Stack overflow detection + UART log & Tests passing \\
  \#344 & BMS SoC thresholds & Warning 25\%, critical 5\% + config & Tests passing \\
  \#345 & Remove event-flag infra & Dead code removal (zero callers) & Tests passing \\
  \#346 & ASCII source cleanup & Replace non-ASCII in 254 files + Doxyfile & Tests passing \\
  \#349 & Gateway RESET\_ACK & Missing case in dispatchControlFrame() & Tests passing \\
  \#350 & Gap analysis docs & This document & Open \\
  \#351 & UI dashboard & 7 pages with mock data (shadcn/Tailwind) & Manual test \\
  \bottomrule
\end{tabular}
\end{center}

\newpage

% ============================================================
% SECTION 2: e2-studio-star-rx72n-firmware/
% ============================================================
\section{Directory: \texttt{e2-studio-star-rx72n-firmware/}}

\begin{dirbox}
  \textbf{Language:} C23 (GNURX toolchain) + ThreadX RTOS\\
  \textbf{Build:} CMake (portable) + Renesas e\textsuperscript{2} Studio (IDE)\\
  \textbf{Hardware:} Renesas RX72N (240 MHz, 4 MB Flash, 512 KB SRAM)\\
  \textbf{Overall:} Production-ready. All 9 tasks, all drivers, communication stack complete.
\end{dirbox}

\subsection{What Is Complete}

\subsubsection{ThreadX Tasks (9 of 9 implemented)}

\begin{center}
\small
\begin{tabular}{L{6.5cm}L{4cm}L{1.5cm}L{2cm}}
  \toprule
  \textbf{Task File} & \textbf{Purpose} & \textbf{Lines} & \textbf{Rate} \\
  \midrule
  \filepath{src/tasks/app\_main\_task.c} & RTOS entry, task creation & 538 & startup \\
  \filepath{src/tasks/comm\_task.c} & SPI/USB frame dispatch & 2060 & IRQ-driven \\
  \filepath{src/tasks/motor\_control\_task.c} & 4-motor PID loop & 2870 & 100 Hz \\
  \filepath{src/tasks/telemetry\_task.c} & Sensor aggregation + proto & 1870 & 20 Hz \\
  \filepath{src/tasks/bms\_monitor\_task.c} & Battery fault detection & 1214 & 10 Hz \\
  \filepath{src/tasks/obstacle\_detect\_task.c} & HC-SR04 sensor fusion & 2202 & 50 Hz \\
  \filepath{src/tasks/temp\_sensor\_task.c} & DS18B20 1-Wire polling & 1227 & 1 Hz \\
  \filepath{src/tasks/led\_status\_task.c} & Status LED patterns & 411 & 20 Hz \\
  \filepath{src/tasks/watchdog\_monitor\_task.c} & IWDT refresh + stack check & 442 & 100 Hz \\
  \bottomrule
\end{tabular}
\end{center}

\subsubsection{Communication Stack}

\begin{center}
\small
\begin{tabular}{L{4.5cm}L{1.5cm}L{8cm}}
  \toprule
  \textbf{Library} & \textbf{Lines} & \textbf{Capability} \\
  \midrule
  \filepath{libs/rx\_spi\_comm/} & 2570 & RSPI2 @ 10 Mbps, zero-copy buffers \\
  \filepath{libs/rx\_usb\_comm/} & 1575 & USB CDC-ACM, 3 virtual COM ports \\
  \filepath{libs/rx\_comm\_manager/} & 2056 & Multi-transport coordinator \\
  \filepath{libs/rx\_fec/} & 869 & Convolutional FEC, K=7, rate 1/2 \\
  \filepath{libs/rx\_harq/} & 675 & Chase combining, configurable retries \\
  \filepath{libs/rx\_frame/} & 895 & SYNC+SEQ+LEN+TYPE+FLAGS+CRC-32 \\
  \filepath{libs/rx\_crc/} & 2923 & CRC-32 (hardware) + CRC-8 \\
  \filepath{libs/rx\_session/} & 544 & Cross-transport sequence continuity \\
  \bottomrule
\end{tabular}
\end{center}

\subsubsection{Hardware Drivers}

\begin{center}
\small
\begin{tabular}{L{4cm}L{1.5cm}L{8.5cm}}
  \toprule
  \textbf{Driver} & \textbf{Lines} & \textbf{Hardware} \\
  \midrule
  \filepath{libs/rx\_motor/} & 1936 & DRV8243S H-bridge, PWM dead-time insertion \\
  \filepath{libs/rx\_pid/} & 1582 & Discrete-time PID with anti-windup \\
  \filepath{libs/rx\_encoder/} & 1969 & MTU1/2 + TPU quadrature encoder channels \\
  \filepath{libs/rx\_drv8243/} & 1983 & DRV8243S SPI configuration driver \\
  \filepath{libs/rx\_bq4050/} & 1293 & BQ4050 BMS via SMBus (RIIC1) \\
  \filepath{libs/rx\_ds18b20/} & 1429 & DS18B20 temperature via 1-Wire \\
  \filepath{libs/rx\_hcsr04/} & 4761 & HC-SR04 ultrasonic, IRQ-based echo \\
  \filepath{libs/rx\_hal/} & 16043 & Full RX72N HAL (GPIO/SPI/UART/ADC/PWM) \\
  \filepath{libs/rx\_bus/} & 11941 & Bus manager (I2C/SPI/1-Wire/UART/SMBus) \\
  \bottomrule
\end{tabular}
\end{center}

\subsubsection{WireMessage Command Dispatcher -- Confirmed Complete}

\begin{okbox}
  \filepath{comm\_task.c:1959} -- \texttt{internal\_handle\_command\_frame()} handles all incoming commands.
\end{okbox}

\begin{lstlisting}[language=C, caption={Command dispatcher -- handles 4 message types}]
/* SetVelocityRequest -> shared_data_set_motor_command() */
/* EmergencyStopRequest -> shared_data_trigger_estop()  */
/* SetPIDGainsRequest -> shared_data_set_pid_gains()    */
/* SetRetransmitConfigRequest -> rx_comm_manager_set_auto_retransmit() */
\end{lstlisting}

\subsubsection{Active Firmware PRs}

\begin{prbox}
  \textbf{\inprogress{341}} -- Frame decoder resync: bounded linear scan (\texttt{k\_frame\_max\_scan\_bytes=1036}), new public \texttt{rx\_frame\_decode\_with\_resync()} API, 3 new Unity tests.\\[4pt]
  \textbf{\inprogress{342}} -- Telemetry transport failover: telemetry was hardcoded to USB only (\texttt{k\_comm\_channel\_usb}). Adds \texttt{internal\_select\_transport()} with USB-first, SPI-fallback. Critical production fix.\\[4pt]
  \textbf{\inprogress{343}} -- ThreadX stack checking: enables \texttt{TX\_ENABLE\_STACK\_CHECKING}, adds \texttt{rx\_stack\_monitor} module, overflow handler logs thread name then halts (IWDT resets in 2 s). 7 Unity tests.\\[4pt]
  \textbf{\inprogress{344}} -- BMS SoC thresholds: adds \texttt{bms\_monitor\_config\_t} with configurable \texttt{soc\_warning\_pct} (25\%) and \texttt{soc\_critical\_pct} (5\%). 14 new unit tests.\\[4pt]
  \textbf{\inprogress{345}} -- Event-flag cleanup: removes dead \texttt{TX\_EVENT\_FLAGS\_GROUP} infrastructure from \texttt{shared\_data} (\texttt{shared\_data\_wait\_event()} had zero callers). All 48 tests pass.\\[4pt]
  \textbf{\inprogress{346}} -- ASCII source cleanup: replaces all non-ASCII characters across 254 files. Adds \texttt{Doxyfile} and Makefile targets (\texttt{doxygen}, \texttt{doxygen-pdf}).
\end{prbox}

\subsection{What Is Missing}

\subsubsection{Gap F-1: NVS Configuration Persistence (HIGH)}

\begin{warnbox}
  \textbf{Impact:} PID gains received via \texttt{SetPIDGainsRequest} are applied in RAM via
  \texttt{shared\_data\_set\_pid\_gains()} but NOT saved to flash. Every reboot resets to
  compile-time defaults -- any tuning session must be repeated.\\
  \textbf{Hardware:} RX72N has 64 KB Data Flash (separate from code flash), ideal for NVS.
\end{warnbox}

\begin{lstlisting}[caption={Files to create for NVS}]
e2-studio-star-rx72n-firmware/
  libs/rx_nvs/
    inc/rx_nvs.h    -- Key-value store API
    src/rx_nvs.c    -- Data Flash implementation
  tests/
    test_rx_nvs.c   -- Unity tests (mock flash)
\end{lstlisting}

\begin{lstlisting}[language=C, caption={NVS API (rx\_nvs.h)}]
/** @brief Write named config blob to Data Flash (max 4096 bytes). */
rx_err_t rx_nvs_write(const char* key,
                       const void* data, uint16_t len);
/** @brief Read named config blob from Data Flash. */
rx_err_t rx_nvs_read(const char* key,
                      void* data, uint16_t* len);
/** @brief Factory reset: erase all Data Flash. */
rx_err_t rx_nvs_erase_all(void);
\end{lstlisting}

\textbf{Integration points:}

\begin{lstlisting}
// On PID update (comm_task.c):
shared_data_set_pid_gains(...);
rx_nvs_write("pid", &gains, sizeof(gains));

// On boot (main.c, before tx_kernel_enter):
rx_nvs_read("pid", &gains, &len);
shared_data_set_pid_gains(&gains);
\end{lstlisting}

\textbf{Estimated effort:} 4--6 hours.

\subsubsection{Gap F-2: OTA Firmware Update Handler (LOW -- Future Enhancement)}

\begin{warnbox}
  \textbf{Scope note:} OTA is not required for the current prototype phase. Hardware can be
  reflashed over JTAG/SWD during development. This is a future enhancement for production deployment.\\
  \textbf{Gateway:} All 10 \texttt{FirmwareUpdateService} RPCs return \texttt{codes.Unimplemented}.\\
  \textbf{Effort when needed:} 2--3 days (firmware state machine + gateway implementation).
\end{warnbox}

\newpage

% ============================================================
% SECTION 3: star-gateway/
% ============================================================
\section{Directory: \texttt{star-gateway/}}

\begin{dirbox}
  \textbf{Language:} Go 1.22+\\
  \textbf{Build:} \texttt{go build ./cmd/star-gateway}\\
  \textbf{Role:} Bridges gRPC clients (UI, ROS2) to RX72N via SPI/USB CDC\\
  \textbf{Test coverage:} 80.8\% -- enforced in CI via \texttt{gateway.yml}\\
  \textbf{Overall:} Transport, framing, FEC, HARQ production-ready. FirmwareUpdateService is a future enhancement.
\end{dirbox}

\subsection{Directory Structure}

\begin{lstlisting}[caption={star-gateway/ layout}]
star-gateway/
  cmd/
    star-gateway/main.go         -- Entry point
    virtual_rx72n/               -- Hardware simulator (95% complete)
  internal/
    transport/                   -- Layer 1: SPI + USB CDC transports
    frame/                       -- Layer 2: SYNC+SEQ+LEN+TYPE+FLAGS+CRC-32
    fec/                         -- FEC: Viterbi K=7, rate 1/2
    harq/                        -- HARQ: Chase Combining retransmit
    service/
      motor_control.go           -- MotorControlService (90%)
      telemetry.go               -- TelemetryService (90%)
      battery.go                 -- BatteryManagementService (100%)
      configuration.go           -- ConfigurationService (85%)
      firmware.go                -- FirmwareUpdateService (0% -- future)
      gateway_service.go         -- GatewayService for ROS2 bridge (85%)
    manager/                     -- Transport manager, RESET_ACK dispatch
    link/                        -- SPI link layer
\end{lstlisting}

\subsection{Test Coverage by Package}

\begin{center}
\small
\begin{tabular}{L{5cm}L{2.5cm}L{6cm}}
  \toprule
  \textbf{Package} & \textbf{Coverage} & \textbf{Status} \\
  \midrule
  \texttt{internal/controller} & 100\% & \done \\
  \texttt{internal/fec} & 100\% & \done \\
  \texttt{internal/frame} & 100\% & \done \\
  \texttt{internal/transport} & 92.1\% & \done \\
  \texttt{internal/manager} & 97.4\% & \done \\
  \texttt{internal/service} & 94.4\% & \partialstatus\ (firmware.go is future) \\
  \texttt{internal/harq} & 87.0\% & \done \\
  \texttt{internal/link} & 83.3\% & \done \\
  \bottomrule
\end{tabular}
\end{center}

\subsection{Virtual RX72N Simulator (\texttt{cmd/virtual\_rx72n/})}

\begin{center}
\small
\begin{tabular}{L{9cm}L{5cm}}
  \toprule
  \textbf{Feature} & \textbf{Status} \\
  \midrule
  Full SPI frame protocol & \done \\
  PING/PONG heartbeat with counter echo & \done \\
  RESET/RESET\_ACK session sync & \done \\
  VelocityCommand + EmergencyStop handling & \done \\
  Simulated TelemetryData (static values) & \done \\
  Sequence number tracking with wraparound & \done \\
  ACK frame generation (FlagRequiresAck) & \done \\
  Configurable latency (SIM\_LATENCY\_MS env var) & \done \\
  Graceful shutdown with socket cleanup & \done \\
  Motor dynamics simulation G(s)=3.665/(0.075s+1) & \missing \\
  Fault injection API (CRC errors, packet loss) & \missing \\
  BMS state-of-charge discharge simulation & \missing \\
  \bottomrule
\end{tabular}
\end{center}

\textbf{To enable integration testing without hardware}, enhance the simulator with motor dynamics:

\begin{lstlisting}[language=Go, caption={Motor dynamics model for virtual\_rx72n/motor\_sim.go}]
// First-order: G(s) = K/(tau*s + 1),  K=3.665 rad/s/V, tau=0.075 s
// Discrete backward Euler at T=0.010 s:
//   v[k] = 0.1248*K*u[k] + 0.8752*v[k-1]
alpha := math.Exp(-motorSamplePeriod / motorTimeConstant)
m.velocityRadps = (1-alpha)*motorGain*inputVoltage + alpha*m.velocityRadps
m.positionRad  += m.velocityRadps * motorSamplePeriod
ticks          := int64(m.positionRad / (2*math.Pi) * ticksPerRev)
m.encoderTicks.Store(ticks)
\end{lstlisting}

\subsection{What Is Missing}

\subsubsection{Gap G-1: RESET\_ACK Auto-Response (IN PROGRESS)}

\begin{prbox}
  \textbf{\inprogress{349}} -- \texttt{dispatchControlFrame()} was missing \texttt{case frame.FrameTypeReset:}, causing firmware-initiated RESET frames to be silently discarded. Adds \texttt{sendResetAck()} modelled after \texttt{sendPong()} with operations-gate bypass to avoid deadlock. All 13 packages pass.
\end{prbox}

\subsubsection{Gap G-2: E-Stop Priority Queue (MEDIUM)}

\begin{warnbox}
  \textbf{File:} \filepath{internal/service/motor\_control.go} -- TODO comment: ``Implement priority framing for E-Stop.''\\
  \textbf{Impact:} During link congestion, E-Stop may be delayed behind pending ACKs.\\
  \textbf{Required:} Bypass the normal HARQ queue for emergency stop frames.\\
  \textbf{Estimated effort:} 2--3 hours.
\end{warnbox}

\subsubsection{Gap G-3: SetMotorPower Not Dispatching (MEDIUM)}

\texttt{SetMotorPower} RPC exists but does not send \texttt{MotorPowerCommand} to the RX72N. The firmware dispatcher also has no handler for it. Requires a coordinated change to \texttt{wire.proto}, firmware dispatcher, and gateway service.

\textbf{Estimated effort:} 3--4 hours (proto + firmware + gateway).

\subsubsection{Gap G-4: Virtual RX72N Motor Dynamics (LOW)}

Simulator returns static telemetry regardless of velocity commands. Adding motor dynamics (see code above) would enable odometry and PID testing without hardware.

\textbf{Estimated effort:} 1--2 days.

\newpage

% ============================================================
% SECTION 4: star-ros2/
% ============================================================
\section{Directory: \texttt{star-ros2/}}

\begin{dirbox}
  \textbf{Language:} C++17, ROS2 Jazzy\\
  \textbf{Build:} \texttt{colcon build} + \texttt{clang-format} enforcement\\
  \textbf{Role:} Middleware -- bridges gateway gRPC to robot topics, manages safety, hosts RPLiDAR\\
  \textbf{Overall:} Core bridge nodes complete. Hardware integration untested. SLAM stack not started.
\end{dirbox}

\subsection{Package Structure}

\begin{lstlisting}[caption={star-ros2/ package layout}]
star-ros2/
  star_gateway_bridge/       -- gRPC client <-> ROS2 topics (COMPLETE)
  star_spi_bridge/           -- SPI driver ROS2 node (code complete, untested)
  star_safety_monitor/       -- Safety watchdog (40% test coverage)
  star_bringup/              -- Launch files (EMPTY SKELETON)
  star_nav/                  -- Navigation stack (0%)
\end{lstlisting}

\subsection{What Is Complete}

\subsubsection{\texttt{star\_gateway\_bridge} (100\%)}

\begin{okbox}
  Fully implemented. Bridges gateway gRPC service to ROS2 topics.
\end{okbox}

\begin{center}
\small
\begin{tabular}{L{5cm}L{4cm}L{5cm}}
  \toprule
  \textbf{Function} & \textbf{ROS2 Interface} & \textbf{gRPC Call} \\
  \midrule
  Teleop forwarding & \texttt{/velocity\_command} (sub) & \texttt{GetTeleopCommand()} \\
  Telemetry bridge & \texttt{/robot\_status} (pub) & \texttt{ForwardTelemetry()} \\
  PID gain update & \texttt{/pid\_gains} (sub) & \texttt{SetPIDGains()} \\
  \bottomrule
\end{tabular}
\end{center}

\subsubsection{\texttt{star\_spi\_bridge} (85\% -- code complete, untested)}

\begin{warnbox}
  All source code is written. This package has \textbf{never been run on real hardware} (RPi5 + RX72N).
  This is the final integration point before the robot can physically move.
\end{warnbox}

\textbf{Known safety gap in} \filepath{src/star\_spi\_driver\_node.cpp}:

\begin{lstlisting}[language=C++, caption={on\_deactivate() must send zero velocity before closing SPI}]
// CURRENT (unsafe): closes SPI while motors may still be spinning
CallbackReturn on_deactivate(const State&) override {
    transport_->close();
    return CallbackReturn::SUCCESS;
}

// REQUIRED (safe):
CallbackReturn on_deactivate(const State&) override {
    auto zero = build_velocity_command(0.0f, 0.0f, 0.0f, 0.0f);
    transport_->send(zero);
    std::this_thread::sleep_for(std::chrono::milliseconds(50));
    transport_->close();
    return CallbackReturn::SUCCESS;
}
\end{lstlisting}

\subsection{What Is Missing}

\subsubsection{Gap R-1: SPI Bridge Hardware Testing (HIGH)}

\begin{criticalbox}
  \textbf{Effort:} 4--8 hours on RPi5 + RX72N hardware.\\
  \textbf{Validates:} Frame encoding, HARQ retransmit, telemetry reception, encoder feedback,
  100 Hz polling requirement (must not miss cycles to avoid 500 ms RX72N timeout).
\end{criticalbox}

\textbf{Pre-test checklist:}
\begin{itemize}[noitemsep]
  \item Fix \texttt{on\_deactivate()} zero-velocity safety (above)
  \item Verify SPI CS polarity (active-low) matches firmware
  \item Confirm SPI Mode 0 @ 10 MHz in RPi5 device tree
  \item Verify COPI/CIPO pin assignments against hardware pinout
\end{itemize}

\subsubsection{Gap R-2: \texttt{star\_bringup} Launch Files (HIGH)}

\begin{criticalbox}
  \filepath{star\_bringup/launch/} is completely empty. Cannot start the system without
  manually launching each node.
\end{criticalbox}

\begin{lstlisting}[caption={Required launch files}]
star_bringup/launch/
  star_robot.launch.py        -- Full system (bridge + SPI + safety + nav)
  star_simulation.launch.py   -- Simulator (Gazebo + virtual_rx72n)
  star_teleop.launch.py       -- Manual control only
  star_nav.launch.py          -- Autonomous (adds SLAM + NavStack)
\end{lstlisting}

\textbf{Estimated effort:} 4--6 hours.

\subsubsection{Gap R-3: Safety Monitor Test Coverage (MEDIUM)}

\begin{warnbox}
  \filepath{star\_safety\_monitor/} -- 3 of 10 test cases use \texttt{GTEST\_SKIP}.\\
  Battery critically low, over-temperature, and communication timeout paths are not tested.\\
  \textbf{Estimated effort:} 4--6 hours.
\end{warnbox}

\subsubsection{Gap R-4: RPLiDAR C1 ROS2 Integration (MEDIUM)}

\begin{warnbox}
  \textbf{Location note:} The RPLiDAR C1 sensor connects to the \textbf{Raspberry Pi 5} (not the RX72N).
  It runs as a ROS2 node on the host, not as firmware. Use the existing
  \texttt{rplidar\_ros2} community package.
\end{warnbox}

\begin{center}
\small
\begin{tabular}{L{4cm}L{10cm}}
  \toprule
  \textbf{Item} & \textbf{What Is Needed} \\
  \midrule
  ROS2 package & Add \texttt{rplidar\_ros2} to workspace dependencies \\
  Launch integration & Include lidar node in \texttt{star\_bringup} launch files \\
  Topic remapping & \texttt{/scan} topic available for SLAM consumption \\
  Hardware mount & Physical UART/USB connection from RPi5 to RPLiDAR \\
  \bottomrule
\end{tabular}
\end{center}

\textbf{Estimated effort:} 4--8 hours (driver setup + launch integration + SLAM config).

\subsubsection{Gap R-5: SLAM Stack (LOW -- requires RPLiDAR + IMU first)}

\begin{center}
\small
\begin{tabular}{L{5cm}L{9cm}}
  \toprule
  \textbf{Component} & \textbf{Status} \\
  \midrule
  RTAB-Map (mapping + localization) & Not installed or configured \\
  \texttt{robot\_localization} EKF & Not configured \\
  Nav2 (path planning) & Not installed or configured \\
  IMU (MPU-6050 or ICM-42688) & Hardware not yet chosen -- connects to RPi5 \\
  \bottomrule
\end{tabular}
\end{center}

\textbf{Prerequisite:} RPLiDAR C1 ROS2 integration (Gap R-4) must work first.

\newpage

% ============================================================
% SECTION 5: star-ui/
% ============================================================
\section{Directory: \texttt{star-ui/}}

\begin{dirbox}
  \textbf{Language:} TypeScript, React 19, Tailwind v4, shadcn/ui\\
  \textbf{Build:} \texttt{npm run build} (Vite)\\
  \textbf{Overall:} All 7 pages scaffolded with mock data in PR \#351. Real data connection is the next step.
\end{dirbox}

\subsection{In Progress -- PR \#351}

\begin{prbox}
  \textbf{\inprogress{351}} -- Complete STAR dashboard frontend. Adds 7 pages with dark-themed
  shadcn/ui components and a \texttt{MockGatewayClient} that returns realistic simulated data.
  The UI is fully navigable and interactive without hardware.\\[4pt]
  \textbf{Pages added:} Dashboard, Controller, Telemetry, Battery, Motors, Configuration, Firmware\\[4pt]
  \textbf{Architecture:} \texttt{GatewayClient} interface + \texttt{MockGatewayClient} implementation.
  Switching to real data requires implementing the real \texttt{GatewayClient} class.\\[4pt]
  \textbf{Status:} Manual test plan only (no automated tests yet).
\end{prbox}

\subsection{What Is Missing}

\subsubsection{Gap U-1: Real Gateway Connection (HIGH -- replaces MockGatewayClient)}

\begin{criticalbox}
  \textbf{Impact:} The UI is fully built (PR \#351) but uses mock data. Connecting to the real
  gateway requires a gRPC-Web or REST bridge between the browser and Go gRPC server.
\end{criticalbox}

\textbf{Options:}
\begin{enumerate}[noitemsep]
  \item \textbf{grpc-gateway} (recommended, ~4 hrs): REST/JSON transcoding to existing gRPC services. No extra process.
  \item \textbf{Envoy proxy} (production-grade): HTTP/2 proxy as a sidecar on RPi5. More complex.
  \item \textbf{connect-go} (modern): Replace grpc-go with connect-go for native HTTP/1.1.
\end{enumerate}

\begin{lstlisting}[language=Go, caption={grpc-gateway setup in cmd/star-gateway/main.go}]
mux := runtime.NewServeMux()
opts := []grpc.DialOption{grpc.WithTransportCredentials(
    insecure.NewCredentials())}
// Register all 5 services for REST transcoding:
starv1.RegisterMotorControlServiceHandlerFromEndpoint(
    ctx, mux, "localhost:50051", opts)
starv1.RegisterTelemetryServiceHandlerFromEndpoint(
    ctx, mux, "localhost:50051", opts)
// ... repeat for Battery, Configuration, Firmware
http.ListenAndServe(":8080", mux)
\end{lstlisting}

\textbf{Estimated effort:} 4--8 hours (bridge setup + real GatewayClient class in UI).

\subsubsection{Gap U-2: Emergency Stop Button (HIGH -- safety critical)}

\begin{criticalbox}
  No E-Stop button exists in the UI. Operators can only stop via gamepad button --
  no visible on-screen control is present. This is a safety gap.
\end{criticalbox}

\textbf{Required:} Prominent red ``EMERGENCY STOP'' button always visible in the header. Keyboard shortcut (Space or Escape). Calls the E-Stop RPC. Shows ACTIVE/CLEARED state.

\textbf{Estimated effort:} 2--3 hours (after real gateway connection working).

\subsubsection{Gap U-3: Tests (MEDIUM)}

Currently 0\% test coverage. Vitest is configured.

\textbf{Recommended:}
\begin{itemize}[noitemsep]
  \item Vitest + React Testing Library for unit tests (component behaviour)
  \item Playwright for E2E (full page interaction)
\end{itemize}

\textbf{Estimated effort:} 8--12 hours for basic coverage.

\newpage

% ============================================================
% SECTION 6: star-proto/
% ============================================================
\section{Directory: \texttt{star-proto/}}

\begin{dirbox}
  \textbf{Language:} Protocol Buffers 3 (buf toolchain)\\
  \textbf{Targets:} Go, TypeScript, nanopb C, C++ (gRPC)\\
  \textbf{Overall:} Schemas are complete. One known nanopb options bug. All generated code is checked in.
\end{dirbox}

\subsection{Schema Inventory}

\begin{center}
\small
\begin{tabular}{L{5.5cm}L{2cm}L{2.5cm}L{4cm}}
  \toprule
  \textbf{Proto File} & \textbf{Msgs} & \textbf{Services} & \textbf{Status} \\
  \midrule
  \filepath{common.proto} & 5 & -- & \done \\
  \filepath{motor\_control.proto} & 9 & 1 (5 RPC) & \done \\
  \filepath{telemetry.proto} & 7 & 1 (3 RPC) & \done \\
  \filepath{battery\_management.proto} & 15 & 1 (10 RPC) & \done \\
  \filepath{configuration.proto} & 13 & 1 (7 RPC) & \done \\
  \filepath{firmware\_update.proto} & 12 & 1 (10 RPC) & \done\ (future use) \\
  \filepath{gateway\_service.proto} & 6 & 1 (3 RPC) & \done \\
  \filepath{diagnostics.proto} & 2 & -- & \done \\
  \filepath{controller.proto} & 1 & -- & \done \\
  \filepath{wire.proto} & 2 & -- & \done \\
  \midrule
  \textbf{Total} & \textbf{72} & \textbf{5 svcs / 38 RPC} & \\
  \bottomrule
\end{tabular}
\end{center}

\subsection{Known Bug: nanopb motor\_status Truncation}

\begin{warnbox}
  \textbf{File:} \filepath{nanopb/star/v1/motor\_control.options}\\
  \textbf{Bug:} \texttt{star.v1.SetVelocityResponse.motor\_status max\_count:2}\\
  \textbf{Fix:} Change to \texttt{max\_count:4} (the robot has 4 motors: FL, FR, BL, BR)\\
  \textbf{Impact:} Response messages silently truncate motors 3 and 4.\\
  \textbf{Fix effort:} 30 minutes -- edit file, run \texttt{make proto-gen \&\& make proto-gen-firmware}.
\end{warnbox}

\subsection{Missing Wire Protocol Messages (for OTA -- future)}

{\sloppy When OTA is implemented, \filepath{wire.proto} will need new types added
to the \texttt{WireMessage} oneof (not required for the current prototype):}

\begin{itemize}[noitemsep,topsep=4pt]
  \item \texttt{FirmwareChunk} -- binary chunk payload
  \item \texttt{BeginUpdateRequest} -- OTA session initialiser
  \item \texttt{FinalizeUpdateRequest} -- CRC validation trigger
\end{itemize}

\newpage

% ============================================================
% SECTION 7: CI/CD (.github/workflows/)
% ============================================================
\section{Directory: \texttt{.github/workflows/}}

\begin{dirbox}
  \textbf{Existing CI:} proto.yml, gateway.yml, ros2.yml (all working)\\
  \textbf{Missing CI:} UI build, docs compile\\
  \textbf{Note:} Firmware CI is \textbf{not planned} -- the GNURX cross-compiler toolchain is too large to install in CI runners.
\end{dirbox}

\subsection{Existing Workflows}

\begin{center}
\small
\begin{tabular}{L{3cm}L{4cm}L{7cm}}
  \toprule
  \textbf{Workflow} & \textbf{Triggers} & \textbf{Steps} \\
  \midrule
  \texttt{proto.yml} & Push to \texttt{star-proto/} & buf lint, breaking check, generate, Go/TS/nanopb tests \\
  \texttt{gateway.yml} & Push to \texttt{star-gateway/} & build, test (-race), lint, security scan, cross-compile (amd64/arm64) \\
  \texttt{ros2.yml} & Push to \texttt{star-ros2/} & colcon build, ament\_lint, code review \\
  \texttt{proto-gen.yml} & Push to \texttt{*.proto} & buf format, generate, check for uncommitted changes \\
  \bottomrule
\end{tabular}
\end{center}

\subsection{What Is Missing}

\subsubsection{Gap C-1: UI Build CI (\texttt{ui.yml}) -- MEDIUM}

\begin{warnbox}
  No CI exists for \filepath{star-ui/}. TypeScript errors and broken builds can merge undetected.
\end{warnbox}

\begin{lstlisting}[caption={Required steps for ui.yml}]
steps:
  - npm ci
  - npm run type-check   # TypeScript strict mode
  - npm run lint         # ESLint
  - npm run test         # Vitest unit tests
  - npm run build        # Production Vite bundle (catches issues)
\end{lstlisting}

\textbf{Estimated effort:} 2--3 hours.

\subsubsection{Gap C-2: Documentation Build CI (\texttt{docs.yml}) -- LOW}

\begin{lstlisting}[caption={Required steps for docs.yml}]
steps:
  - apt-get install texlive-xetex texlive-fonts-extra
  - cd docs && make        # Compile star_documentation.pdf
  - Upload PDF as artifact
\end{lstlisting}

\textbf{Estimated effort:} 1--2 hours.

\newpage

% ============================================================
% SECTION 8: IMPLEMENTATION ROADMAP
% ============================================================
\section{Implementation Roadmap}

\subsection{Phase 0: Merge Open PRs (in flight)}

All 7 open firmware/infra PRs are ready (tests passing). Merge in order to avoid conflicts:

\begin{enumerate}[noitemsep]
  \item \inprogress{345} -- Remove event-flag infra (no conflicts, merge first)
  \item \inprogress{342} -- Telemetry USB/SPI failover (critical production fix)
  \item \inprogress{341} -- Frame decoder resync
  \item \inprogress{343} -- ThreadX stack checking
  \item \inprogress{344} -- BMS SoC thresholds (depends on \#345 merged first)
  \item \inprogress{346} -- ASCII source cleanup (touches many files, merge last)
  \item \inprogress{349} -- Gateway RESET\_ACK fix
  \item \inprogress{351} -- UI dashboard (review mock architecture)
\end{enumerate}

\subsection{Phase 1: Quick Wins (1--2 hours)}

\begin{enumerate}[noitemsep]
  \item \textbf{Fix nanopb motor\_status bug} (\filepath{star-proto/}): change \texttt{max\_count:2} to \texttt{max\_count:4}
  \item \textbf{Fix SPI bridge zero-velocity safety} (\filepath{star-ros2/star\_spi\_bridge/}): send zero-vel in \texttt{on\_deactivate()}
\end{enumerate}

\subsection{Phase 2: First Robot Motion (2--4 days on hardware)}

\begin{enumerate}[noitemsep]
  \setcounter{enumi}{2}
  \item \textbf{Hardware bring-up:} verify SPI/COPI/CIPO/CS wiring matches\\firmware pinout
  \item \textbf{Run SPI bridge on hardware:} test with virtual\_rx72n first, then real firmware
  \item \textbf{Create star\_bringup launch files:} minimum \texttt{star\_robot.launch.py}
  \item \textbf{Validate teleop path:} gamepad $\to$ gateway $\to$ ROS2 $\to$ SPI bridge $\to$ motors spin
\end{enumerate}

\subsection{Phase 3: Real UI Data (1--2 weeks)}

\begin{enumerate}[noitemsep]
  \setcounter{enumi}{6}
  \item \textbf{Implement real GatewayClient} in \filepath{star-ui/} (replaces MockGatewayClient)
  \item \textbf{Add grpc-gateway} or REST bridge to gateway (browser access to gRPC services)
  \item \textbf{Emergency Stop button} (safety-critical UI component)
  \item \textbf{UI CI} (\texttt{ui.yml} workflow)
\end{enumerate}

\subsection{Phase 4: Robustness (1--2 weeks)}

\begin{enumerate}[noitemsep]
  \setcounter{enumi}{10}
  \item \textbf{NVS library} (\filepath{libs/rx\_nvs/}): PID gain persistence across reboot
  \item \textbf{E-Stop priority queue} in gateway: bypass normal HARQ queue
  \item \textbf{Safety monitor tests}: implement 3 skipped test cases
  \item \textbf{Virtual RX72N motor dynamics}: enables sim-based PID tuning
\end{enumerate}

\subsection{Phase 5: SLAM and Autonomy (future)}

\begin{enumerate}[noitemsep]
  \setcounter{enumi}{14}
  \item RPLiDAR C1 ROS2 integration on RPi5 (\texttt{rplidar\_ros2} package)
  \item RTAB-Map configuration for ROS2 Jazzy
  \item \texttt{robot\_localization} EKF with encoder odometry
  \item Nav2 path planner
\end{enumerate}

\subsection{Effort Summary}

\begin{center}
\small
\begin{longtable}{L{6cm}L{2.5cm}L{2cm}L{3.5cm}}
  \toprule
  \textbf{Gap} & \textbf{Severity} & \textbf{Effort} & \textbf{Status} \\
  \midrule
  \endhead

  Fix nanopb max\_count:4 & Critical bug & 30 min & Not started \\
  SPI on\_deactivate() zero-vel & Safety & 1 hr & Not started \\
  Merge open PRs (\#341--\#351) & High & -- & \inprogress{341-349} \\
  star\_bringup launch files & High & 4--6 hrs & Not started \\
  SPI bridge hardware testing & High & 4--8 hrs & Not started \\
  Real GatewayClient (UI) & High & 4--8 hrs & \inprogress{351} (mock) \\
  gRPC-Web / REST bridge & High & 4--8 hrs & Not started \\
  Emergency Stop UI button & High (safety) & 2--3 hrs & Not started \\
  Safety monitor test coverage & Medium & 4--6 hrs & Not started \\
  NVS library (rx\_nvs) & Medium & 4--6 hrs & Not started \\
  PID gain persistence & Medium & 2 hrs & Needs NVS first \\
  E-Stop priority queue & Medium & 2--3 hrs & Not started \\
  UI CI (ui.yml) & Medium & 2--3 hrs & Not started \\
  SetMotorPower dispatch & Medium & 3--4 hrs & Not started \\
  Docs CI (docs.yml) & Low & 1--2 hrs & Not started \\
  Virtual RX72N motor dynamics & Low & 1--2 days & Not started \\
  Virtual RX72N fault injection & Low & 4--8 hrs & Not started \\
  RPLiDAR C1 ROS2 (on RPi5) & Low & 4--8 hrs & Not started \\
  OTA firmware update & Future & 2--3 days & Not required (prototype) \\
  SLAM stack & Future & weeks & Not started \\
  \bottomrule
\end{longtable}
\end{center}

\end{document}
